\chapter{Windows, MinGW-w64, and Wine}
\label{ch:windows-wine}

This book's Direct3D chapter described what the D3D9/D3D11/D3D12 backends prove; this chapter
describes the actual machinery that makes proving it possible without ever touching a
Windows machine --- a genuinely elaborate, carefully engineered dev loop that deserves its
own explanation.

\section{The cross-compilation toolchain}

\texttt{cmake/toolchains/mingw-w64.cmake} is a short, standard CMake cross-toolchain file:
it sets \texttt{CMAKE\_SYSTEM\_NAME} to \texttt{Windows}, locates the
\texttt{x86\_64-w64-mingw32-gcc} / \texttt{g++} toolchain (falling back to a 32-bit triple if
the 64-bit one is not found), and isolates the cross-root so CMake's own
\texttt{find\_package} calls never accidentally pick up a host Linux library. The invocation
is exactly what Chapter~\ref{ch:building-cna} already showed:
\texttt{-DCMAKE\_TOOLCHAIN\_FILE=cmake/toolchains/mingw-w64.cmake}, alongside whichever
backend option is being targeted.

\begin{lstlisting}[style=cnashell]
# Worked example: the exact invocation this toolchain file's own header
# comment documents, run from the cna repo root.
cmake -B build-windows \
      -DCMAKE_TOOLCHAIN_FILE=cmake/toolchains/mingw-w64.cmake \
      -DCNA_GRAPHICS_BACKEND=SDL_RENDERER
cmake --build build-windows
\end{lstlisting}

\section{Three Wine prefixes, never one}

A single Wine prefix is not enough, and the project maintains three, each for a distinct
reason. One prefix carries a genuine Microsoft \texttt{d3dcompiler\_47.dll} (installed via
\texttt{winetricks}), needed because Wine's own \emph{builtin} \texttt{d3dcompiler\_47.dll} is
backed by an incomplete shader compiler --- it fails outright on a real alpha-test ternary
expression from the stock effects, and where it does compile something, the result can be
wildly wrong: a 72-bone skinning shader that should compile to about 2{,}160 bytes came out
of the builtin compiler at 1.4~MB, fully unrolled, because it cannot do relative addressing
at all. A second, 32-bit prefix carries a full XNA~4.0 installation --- real .NET Framework 4,
the real XNA~4.0 redistributable, and a real in-prefix \texttt{csc.exe} --- specifically so
the D3D9 oracle (\S\ref{sec:d3d9-oracle}) can compile and run genuine XNA~4.0 code, not a
substitute for it. A third prefix is the one CNA's own D3D11 test suite runs against day to
day, deliberately left alone rather than experimented on.

\section{The self-verifying gate: catching a silent fallback before it corrupts a result}

Every \texttt{scripts/run-wine-dxvk*.sh} / \texttt{run-wine-vkd3d.sh} runner shares one
critical property: none of them simply trust that DXVK or vkd3d-proton is actually in effect.
Each one captures the target program's own log output and greps it for a translation-layer-specific
marker --- a real \texttt{DXVK: <version>} line, or a real \texttt{vkd3d-proton -
applicationVersion: <version>} line --- and refuses to treat the run as valid otherwise,
exiting with a specific error naming the failure as ``a silent WineD3D fallback rather than a
real DXVK-backed Direct3D~11 run.'' This is not a hypothetical concern: without this gate, a
misconfigured prefix could silently render through Wine's own, much less accurate, built-in
Direct3D implementation instead of the real GPU-accelerated translation layer the whole
verification methodology depends on, and every pixel-test result from that run would be
comparing CNA against the wrong thing without anyone noticing.

The gate itself is not unconditional, and reading \texttt{run-wine-dxvk.sh} directly surfaces
two deliberate, narrowly-scoped escape hatches its own comments explain precisely.
\texttt{CNA\_D3D11\_ALLOW\_WINED3D=1} bypasses it for a genuinely deliberate one-off,
non-DXVK diagnostic run --- reproducing an original vanilla-Wine spike result, for instance ---
and is explicitly documented as not for normal test runs.
\texttt{CNA\_D3D11\_SKIP\_DXVK\_GATE=1} covers a different, equally real case: a binary that
legitimately never creates a D3D11 device at all (pure-function mapping-table checks, for
instance), which would never print a DXVK marker line for entirely unrelated, legitimate
reasons --- without this second escape hatch, the gate would misreport that binary's clean run
as a WineD3D fallback it never actually had.

\begin{lstlisting}[style=cnashell]
# Worked example: a normal, gated run -- fails loudly if DXVK didn't
# actually handle it, rather than silently comparing against WineD3D.
scripts/run-wine-dxvk.sh build-windows/CnaD3D11Tests.exe

# A deliberate, one-off diagnostic run against plain WineD3D instead --
# NOT for normal test runs, per the script's own documented contract.
CNA_D3D11_ALLOW_WINED3D=1 scripts/run-wine-dxvk.sh build-windows/CnaD3D11Tests.exe

# A binary that never creates a D3D11 device at all (pure mapping-table
# checks) -- skips the gate rather than being misreported as a fallback.
CNA_D3D11_SKIP_DXVK_GATE=1 scripts/run-wine-dxvk.sh build-windows/D3DCommonMappingTests.exe
\end{lstlisting}

\section{D3D12's extra wrinkle: Proton, not plain Wine}

Chapter~\ref{ch:direct3d-backends} already described the real, root-caused
architecture mismatch behind D3D12's swap-chain crash under plain Wine: Debian's own system
\texttt{dxgi.dll} cannot hand a command queue to vkd3d-proton's separately overridden
\texttt{d3d12.dll}, because vkd3d-proton's own \texttt{dxgi.dll} is only ABI-compatible with
the rest of its own matched build. The concrete fix is a dedicated script that launches
through Valve's own Proton launcher end to end, which supplies a self-consistent DLL set that
vkd3d-proton's override then layers onto correctly. This script carries a hard-coded safety
guard refusing to ever operate against a real personal Wine prefix --- added after a genuine
incident where an earlier version of the same script accidentally launched against a
developer's own \texttt{\textasciitilde/.wine} prefix and hung for nine and a half hours
under a first-run installation wizard.

\section{The D3D9 oracle, mechanically}
\label{sec:d3d9-oracle}

Chapter~\ref{ch:direct3d-backends}'s account of the D3D9 oracle can now be given its full
mechanical explanation. A tool compiles Microsoft's own, unmodified stock-effect \texttt{.fx}
source files to real Direct3D~9 bytecode via the genuine \texttt{d3dcompiler\_47.dll},
cross-built and run through Wine against the dedicated compiler prefix above; a separate
comparison step diffs the result against FNA's own shipped, Microsoft-compiled bytecode ---
61 of 66 compiled variants matched byte-for-byte, and the remaining five (all
\texttt{PixelLighting} vertex variants) were confirmed, not merely assumed, to differ only by
compiler version, after ruling out every optimization-flag combination as the actual cause.

The pixel-diff oracle itself works by rendering the identical declarative scene description
through two independent paths: a real, Wine-hosted XNA~4.0 program (compiled by the genuine
in-prefix \texttt{csc.exe}, run against the XNA-redistributable prefix, itself running
through DXVK rather than WineD3D --- a deliberate methodological choice, so that a pixel diff
never silently measures a driver difference instead of a real CNA difference) and CNA's own
cross-compiled D3D9 backend, run through the same DXVK-verified gate as everything else in
this chapter. A dedicated, mutation-tested comparison script does the actual pixel diff, at
a caller-specified tolerance --- confirmed, via a deliberately corrupted test image, to
correctly report exactly the number of differing pixels a known mutation introduces, not just
a pass/fail verdict.

\subsection{A real D3D9 API restriction, surfaced only by this dev loop}

Real Direct3D~9 hardware/driver behavior, not a CNA logic bug, is exactly what this whole
elaborate dev loop exists to surface reliably --- and one concrete finding from closing out the
D3D9 backend's own verification work is worth naming as proof it does. DXVK's D3D9
implementation, correctly matching real D3D9's own restriction, refused to create a swap chain
using \texttt{D3DFMT\_A8B8G8R8} --- \cnaclass{SurfaceFormat::Color}'s own natural D3D9
equivalent, and a format D3D9 accepts perfectly well for an ordinary \emph{texture}. The error
DXVK's own \texttt{D3D9DeviceEx::ResetSwapChain} raises is explicit: this exact byte order is
simply not one of the small set of formats D3D9 permits for the \emph{primary back buffer}
specifically, a restriction rooted in real display-hardware compatibility, not an arbitrary API
limitation:

\begin{lstlisting}
// Real D3D9 restriction: A8B8G8R8 is legal for a texture, illegal for the
// swap chain's own back buffer. Substituted only at back-buffer creation --
// every ordinary Texture2D backed by Color still requests A8B8G8R8 directly.
if (requested == D3DFMT_A8B8G8R8) return D3DFMT_A8R8G8B8;
\end{lstlisting}

The fix substitutes the back buffer's own display-compatible sibling format,
\texttt{D3DFMT\_A8R8G8B8} (the identical channels, reordered), specifically and only at swap-chain
creation --- every other \cnaclass{SurfaceFormat::Color}-backed resource, an ordinary
\cnaclass{Texture2D} in particular, keeps requesting \texttt{A8B8G8R8} directly, since that
restriction genuinely does not apply there. The substitution stays correct end to end because
the backend's own back-buffer readback path already handles both byte orders, not because the
two formats happen to be interchangeable in general. This is precisely the class of finding
this chapter's own D3D9 oracle and DXVK-gated dev loop are built to catch: a real, confirmed
Direct3D~9 API restriction, discovered only because the verification pipeline runs against a
genuine, correctly-behaving Direct3D~9 implementation rather than a simulated or assumed one.

\section{fna-reference: a different, lighter-weight comparison, and why it uses Mono instead}

Distinct from the rendering-capable D3D9 oracle above, \texttt{tools/fna-reference/} answers
a narrower question --- do CNA's enums, state-object presets, packed vector types, and
viewport math match FNA's own values exactly --- and does so without any Wine or Windows
involvement at all. It runs the real \texttt{FNA.dll} natively on Linux through Mono, dumping
a fixed set of non-rendering values (enum members, default state-object properties, packed
vector round-trips, viewport projection math) to JSON; a companion C\texttt{++} tool performs
the identical dump from CNA's own real implementation, entirely without constructing a
\cnaclass{GraphicsDevice} at all; and a small, dependency-free comparison script asserts every
FNA-side value has a matching CNA-side value, deliberately one-directional so that a
CNA-only, \texttt{NOXNA}-tagged extension is never flagged as a mismatch. This lighter-weight
tool exists specifically because it needs no live rendering device --- extending it to cover
anything that does (a stock effect's actual rendered output, for instance) was evaluated and
explicitly deferred as a substantially larger undertaking, left for the D3D9 oracle to cover
instead. One real, if narrow, divergence this tool did catch: \cnaclass{IndexElementSize}'s
\texttt{SixteenBits} / \texttt{ThirtyTwoBits} values did not match FNA's own \texttt{0} / \texttt{1}
numbering at all --- confirmed and fixed once this comparison surfaced it.

\section{The one caveat every result in this chapter inherits}

Every measurement this elaborate dev loop produces --- the D3D9 oracle's 0-of-31 divergence
count, the D3D11/D3D12 smoke suites, the \texttt{GraphicsProfile.Reach} / \texttt{HiDef} logic
Chapter~\ref{ch:graphicsdevice} already covers as the one CNA backend where that distinction is
genuinely real --- inherits the same single, explicitly-named qualifier, stated plainly and
project-wide rather than buried in one task's own footnote: under Wine plus DXVK,
\texttt{D3DCAPS9} is \emph{synthesized by DXVK itself}, not reported by an authentic,
period-correct Direct3D~9 driver from Intel, AMD, or NVIDIA circa the real Xbox~360/Windows
XNA~4.0 era. The distinction is precise and worth holding onto exactly as stated: the
\emph{logic} this whole verification methodology proves is real --- real \texttt{D3DCAPS9}
fields, read through real \texttt{IDirect3D9} / \texttt{IDirect3DDevice9} calls, with real
profile-ceiling enforcement and no hardcoded ``pretend'' table standing in anywhere. What is
not yet proven is whether the specific \emph{numbers} those real calls return in this
Wine+DXVK dev loop match what an authentic XNA-era GPU driver would have reported --- because
no such hardware exists anywhere in this development environment to compare against. A result
like ``\texttt{IsProfileSupported(HiDef)} returns \texttt{true}'' in this loop proves the
comparison logic correctly reads and evaluates real capability fields; it does not by itself
prove a genuine HiDef-class GPU was actually asked, since DXVK's own synthesized capability set
is standing in for one.

Exactly one remaining task closes this gap, and this book states its status as plainly as the
project's own tracking does: real Windows hardware verification, covering a real device-lost
event fired by an actual alt-tab/RDP/lock-screen transition (not this dev loop's own simulated
context-loss test hook), real \texttt{D3DCAPS9} values from an authentic period driver, real
exclusive fullscreen behavior, and --- the single most definitive remaining measurement in the
whole methodology --- re-running the D3D9 oracle itself on real Windows against a real XNA~4.0
installation rather than through this chapter's own cross-compiled, Wine-hosted substitute for
one. This task is formally marked \texttt{needs\_human} in the project's own tracking --- not
unfinished engineering, but a verification step that requires physical hardware and a human
hand no amount of further Wine/DXVK tooling can substitute for, the identical gap
Chapter~\ref{ch:verification-methodology}'s own cross-platform account already names plainly:
real Windows hardware verification remains open on all three Direct3D backends. Every specific,
checkable claim this book has made about D3D9 pixel authenticity should be read with this
qualifier attached --- ``indistinguishable from real XNA~4.0 running through the same DXVK
path, on this machine,'' a real and rigorously earned result, but not yet ``authentically
indistinguishable on real period Xbox~360/Windows hardware,'' which remains the one honestly
unclosed claim in this entire chapter's own account of itself.
